%
%-----------------------------------------------------------------------
\chapter{Einleitung}%\label{chap:kap1}
%-----------------------------------------------------------------------
    %-----------------------------------------------------------------------
    \section{Stand der Technik} %\label{chap:sec1}
 Durch die Industrialisierung des 19. Jahrhunderts, insbesondere durch die Entwicklung der Eisenbahn und deren Schadensfälle hat August Wöhler mit Versuchen eine neue Technik zur Lösung von zahlreichen Problemen entwickelt. Er hat durch gebrochene Radsatzwellen an Lokomotiven entdeckt, dass Bauteile auch unterhalb der statischen Festigkeitsgrenze beschädigt werden können. 
Sein Versuch, der den Zusammenhang zwischen Bruchlastspielzahl und Ausschlagspannung herstellt zeigt, dass wiederholte, schwingende Belastungen zum Bruch führen können, auch wenn sie weit unterhalb der statischen Festigkeitsgrenze liegen.
Seitdem spielt die dynamische Werkstoffprüfung eine wichtige Rolle und befindet sich mittlerweile im alltäglichen Gebrauch von Firmen.


    %	
    %-----------------------------------------------------------------------
    \section{Motivation} %\label{chap:sec1}
    %-----------------------------------------------------------------------
Die vorliegende Bachelorarbeit dient der Entwicklung einer MATLAB-Auswertesoftware, die die Belastung der Maschine und die Reaktion des Prüflings bewertet. Der komplette Aufgabenbereich ist interessant und spannend, weil er von der Funktion und Bedienung der Anlage bis zum Kundenauftrag mit zertifizierten Prüfschein reicht.  




